\honest{Truly Part Of You}{Eliezer Yudkowsky}{2007}

I begin with Drew McDermott's paper ``Artificial Intelligence Meets Natural Stupidity.'' AI programs stored facts like ``happiness is a state of mind'' as nodes with suggestive English names. McDermott's test was to rename the nodes with meaningless symbols and see whether you still admired the system. If you delete the suggestive names, I add, ``they don't grow back.''

A person who believes ``light is waves'' because a physicist said so has the same kind of node in their head. As McDermott says, ``The whole problem is getting the hearer to notice what it has been told.'' \nb{In McDermott's paper that sentence is about computer programs passing data to each other, a case the paper contrasts with human speech. The post applies it to a person.} My test: ``How would I regenerate this knowledge if it were deleted from my mind?''

Experience connects a belief to other knowledge, to perception and to action. Someone who has only been told facts about ``beavers'' ``may not be able to recognize a beaver'' when they see one. I cite Davidson: if you believe beavers live in deserts, are pure white and weigh 300 pounds, you have no beliefs about beavers at all. \nb{That example concerns beliefs that are mostly false, not true beliefs taken on trust. The article the post cites, by Richard Rorty, explains it by the view that words get their meaning from shared use among speakers, ``not by being paired off with particular experiences or objects.''} Then Wittgenstein: ``A wheel that can be turned though nothing else moves with it, is not part of the mechanism.'' \nb{In Wittgenstein's passage the idle wheel is a private meaning, in a man who ``uses it as we all do.''}

I have asked myself this question since I first read about AI. Could I re-prove the Pythagorean theorem if I forgot the proofs? ``I think so.'' Would I notice the theorem if I forgot it entirely? That is ``harder to boast,'' but if you handed me a 3-4 right triangle and told me the hypotenuse could be calculated, ``I think I would be able to calculate it.'' \nb{In that version someone else notices the problem for the author, which was the part to be tested. The post describes no case in which the test was run and found a belief wanting.}

A shepherd's apprentice who does not understand the pebble-counting system cannot fix an extra pebble, or see how to improve the system. \nb{The shepherd comes from Yudkowsky's ``The Simple Truth,'' the next essay in the book.} Knowledge whose source you hold grows with you.

``Strive to make yourself the source of every thought worth thinking,'' I conclude. \nb{Six weeks earlier, in ``Cached Thoughts,'' Yudkowsky had written that ``no one can think fast enough to think their own thoughts.'' The post does not say which thoughts to regenerate and which to take on trust.}
